\section{Discussion}\label{sec:discussion}

This section interprets the demonstration, identifies five honest limitations, and outlines the gap between what \texttt{mr-audit} v0.0.1 currently does and what a production-grade audit-trail system in a regulated firm would need.

\subsection{For practitioners deploying ML in regulated settings}

The narrow practical claim is that an SR 11-7-shaped audit trail for an ML inference pipeline does not require a commercial platform, a cloud subscription, or a centralized governance system. A one-line decorator and a context manager produce records that satisfy the field-level requirements of \Cref{tab:requirements}. The records are immutable (chain-tied), replayable (verified to bit-identical match in our demonstration), portable (vendor-neutral storage), and reviewer-friendly (regulator-export bundle verifiable with off-the-shelf tools). The marginal engineering cost of doing this on top of an existing pipeline is roughly two days of integration plus a per-prediction logging cost on the order of $10^{-2}$ seconds (matching the throughput in \Cref{tab:instrumentation}).

A second practical takeaway is the value of recording the \emph{full state} of a prediction event, not just the prediction itself. Many firms today record only the prediction value and a timestamp. The audit-trail value comes from the breadth of the record: when an external reviewer asks ``what code, what data, what hyperparameters, what library versions, what random seed produced this number?'', the answer must be in the log without re-running the pipeline. The \texttt{mr-audit} schema captures all of these by default; the engineering investment is in the integration, not in the retention.

\subsection{For the regulatory framework}

The EU AI Act's Article 12 / Article 19 are operational rules, not principles. They specify \emph{what} must be in the log; they do not specify \emph{how} the log should be implemented. Our demonstration is intended as a partial answer to the latter question: in Python, with $\sim$2,000 lines of code, MIT-licensed, on a single laptop. The implementation does not require a cloud provider, a specific database vendor, or a specific ML framework. The implication for firms that have not yet provisioned for AI Act compliance is that the technical infrastructure is achievable with modest engineering effort, and the cost driver will be process and governance rather than tooling.

For SR 11-7, the parallel observation is that the model risk management lifecycle (development, validation, deployment, ongoing monitoring) is well-served by an audit trail that captures every prediction's full state. The current banking-industry norm of recording only inputs and outputs in a separate database, with model parameters versioned in a separate model registry, requires correlation across multiple systems to reconstruct any historical prediction. A unified audit trail collapses this correlation effort.

\subsection{Limitations}\label{sec:limitations}

Five limitations of the present work deserve careful noting.

\textit{1. The demonstration is retrospective.} The 6,825 records were written today against a previously-fit forecast panel. We did not re-run the seven-model pipeline live and write records in real time. The retrospective form is appropriate for showing that the five use cases work end-to-end at realistic scale, but it does not establish that \texttt{mr-audit} would handle a live high-throughput inference workload (e.g., the millions of predictions per day that some bank ML pipelines generate) without re-engineering. The current SQLite single-writer backend's 22.6 records/s throughput is comfortable for a thousand-prediction-per-day workflow but inadequate for a million-per-day workflow. Higher-throughput storage backends (Parquet append shards, write-optimized databases) are a planned 0.2.0 addition.

\textit{2. ``Drift in 44 of 44 months'' is partly a feature, partly a noise issue.} The KS test on a five-feature window over a multi-year volatility series will register drift in nearly every month relative to a fixed early baseline, because realized-volatility distributions shift across regimes. This is the \emph{intended} behavior for SR 11-7 §V monitoring (you want the model risk team to be told when input distributions move) but it is also a \emph{noisy} signal: a team that is told ``drift detected'' every month will calibrate a high tolerance and miss the drift instances that matter. A production deployment would likely use a rolling baseline (e.g., the prior 90 days) rather than a fixed early baseline, and would define drift relative to the magnitude observed during model development. Our demonstration uses the simplest baseline construction for clarity; the package does not preclude more sophisticated baselines.

\textit{3. Integrity protection is hash-chain, not signed.} Tampering with the chain is detectable by \texttt{verify\_chain}, but the chain itself is not cryptographically signed by an external authority. A sufficiently-sophisticated adversary with write access to the log file could rewrite the entire chain forward from the point of tampering, producing a self-consistent but altered chain. Defences against this attack require external attestation: periodic publication of the latest \texttt{record\_id} to a trusted external store (an internal audit ledger, a public blockchain, or a regulator-controlled repository), or signing of the chain by a hardware security module. \texttt{mr-audit} does not provide either; it provides the substrate on which a firm can layer its preferred attestation strategy.

\textit{4. Feature values are stored as summaries by default.} The default \texttt{feature\_values\_summary} captures per-feature numeric values at prediction time. Capturing raw input vectors (e.g., 1-million-dimensional NLP embeddings) would inflate log sizes intolerably and may capture sensitive customer data. The current default is conservative (small log size, no PII risk) but loses information that a more granular drift analysis would need. Firms that need raw-vector logging can configure \texttt{summarize\_features=False}; the privacy and storage trade-offs are theirs to manage.

\textit{5. The package targets Python ML pipelines.} A bank using R, Java, MATLAB, or a vendor-supplied scoring engine cannot use \texttt{mr-audit} directly. The schema (\Cref{tab:schema}) is language-neutral (canonical JSON), so a port to other languages is straightforward in principle, but no port currently exists.

\subsection{What would change our view}\label{sec:falsification}

The paper makes empirical and methodological claims. We state explicitly what would force us to revise each.

\textit{Claim: \texttt{mr-audit}'s schema satisfies SR 11-7, EU AI Act Articles 12 and 19, NIST AI 100-1, and ISO/IEC 42001 Clauses 6.1.2 / 8.3 / 9.1 simultaneously.} We would revise if (a) a supervisory enforcement action against a deployed mr-audit-instrumented system finds the trail insufficient for one of the four frameworks, identifying the specific clause not satisfied; or (b) a peer-reviewed compliance analysis identifies a regulatory clause we have mis-mapped in \Cref{tab:multi_regulation_matrix}. Either revision is incremental, the gap can be added to the schema in a backward-compatible way.

\textit{Claim: hash-chain integrity detects single-record corruption, chain-link tampering, and middle-record deletion (\Cref{tab:adversarial} scenarios A, B, D).} We would revise if an independent reproduction of the adversarial experiment on a different log (different size, different model, different storage backend) yields a non-detection in any of those three scenarios. This is unlikely given the mathematical structure of SHA-256, but the empirical claim is sample-conditional and should be tested against deployment-realistic logs.

\textit{Claim: tail truncation is undetectable by chain verification alone (\Cref{tab:adversarial} scenario C).} We would revise if a method is identified that detects tail truncation from the log alone without external anchoring. We do not believe such a method exists for a self-consistent chain, but we would adopt it if it is.

\textit{Claim: input-vs-output drift comparison enables a covariate-vs-concept-drift distinction (\Cref{tab:drift_interpretation}).} We would revise if a published deployment finds that the four-quadrant framework misclassifies a real drift episode, identifying the case and the misclassification mechanism. The framework is well-established in model-risk practice (e.g., \citealp{KolmogorovSmirnov, PSI_Karakoulas}); the contribution here is the pairing of the two detectors in one package and the audit-log substrate for retrospective evaluation.

\textit{Claim: low-friction integration cost.} We would revise if a deployed pipeline finds the integration cost materially higher than ``one decorator plus one context manager'' (the demonstrated cost on the Khan 2026 panel). Specifically: if a non-trivial bank ML pipeline (say, a credit-scoring engine running 100k predictions per day with multi-model ensembling and human-in-the-loop review) requires more than one engineer-week to integrate, we would document the additional integration patterns required and revise the ``low-friction'' claim accordingly.

The falsification criteria are not asks for permission; they are commitments to update when evidence arrives. We invite that evidence.

\subsection{Roadmap beyond v0.1.0}\label{sec:future}

The v0.1.0 release ships an output-drift module (\texttt{drift/output.py}), symmetric to the input-drift module: KS two-sample plus PSI applied to \texttt{Prediction.value} across windows. Three further additions are planned for v0.2.0 and will not change the audit-record schema, so logs written under the current version remain readable:

\begin{enumerate}[leftmargin=*]
    \item \textbf{Higher-throughput backend.} Parquet-append shard backend with optional compaction, targeting 1{,}000+ records/s on commodity hardware.
    \item \textbf{External attestation hooks.} Plugin interface for periodically publishing the latest \texttt{record\_id} to an external trusted store (Sigstore, an internal HSM, or a blockchain anchor) so that the chain's integrity can be re-verified by an independent party at later dates.
    \item \textbf{Model artifact hashing.} Optional automatic SHA-256 of the underlying model object (a pickled scikit-learn estimator or a PyTorch state-dict) to populate \texttt{model.model\_artifact\_hash}.
\end{enumerate}
